% Conversão fiel da participação como cliente Hoortech; origem imutável RELATORIO_ATUAL.odt.
\hypertarget{__RefHeading___Toc4470_3118176633}{}

% SOURCE text-h-0028
\refstepcounter{subsection}
\subsection*{6.2.2 Sprint 0}
\label{cliente-hoortech:text-h-0028}
\addcontentsline{toc}{subsection}{6.2.2 Sprint 0}

% SOURCE text-p-0712
\originaltab{}Durante a entrega da Sprint 0, tive uma experiência bastante positiva. Os AGES 4 conduziram a apresentação de forma organizada, utilizando um PowerPoint para demonstrar o andamento do cronograma, com as datas dos próximos encontros e as \textit{User Stories }planejadas para cada sprint. Além disso, eles mostraram os primeiros avanços da equipe, incluindo as telas de broadcast, que apresentavam a imagem do usuário junto a uma barra lateral com o diálogo da conversação e várias opções de configurações de texto. Esses protótipos, desenvolvidos no Figma, possuíam alta qualidade visual, o que me deixou bastante satisfeito.\par

% SOURCE text-p-0713
Com base em minha experiência prévia como desenvolvedor em outras edições da AGES, eu me esforcei para oferecer \textit{feedback} construtivo, algo que já notei ser uma dificuldade para outros clientes, que muitas vezes aceitam o trabalho sem oferecer críticas. Isso pode ser prejudicial para os desenvolvedores, que ficam sem parâmetros claros de qualidade. Assim, sugeri alguns ajustes, como o aumento do contraste da caixa de diálogo e o aumento do tamanho das legendas dos botões, os quais foram bem recebidos tanto pelos alunos quanto pelo professor orientador.\par

% SOURCE text-p-0714
Um ponto a ser aprimorado foi o ritmo da apresentação, que aconteceu de forma muito rápida. Esse é um aspecto que já foi comentado pelo meu orientador em edições anteriores da AGES: é importante que os desenvolvedores expliquem o trabalho com mais calma, dando tempo para o cliente analisar as telas com atenção, já que nós, clientes, estamos vendo o produto pela primeira vez. Tive que pedir que voltassem para algumas telas ao final para poder analisá-las com mais calma. Apesar disso, a comunicação da equipe foi excelente, tanto durante a apresentação quanto fora dela, com a equipe fornecendo todas as informações necessárias sobre datas, entregas e compartilhando o material utilizado na apresentação.\par

